\section{The full-unit contact defect and its arithmetic trace derivative}
\label{sec:contact-trace-variation}

The supplied \emph{Regular-inverse contact defects and their arithmetic
return}, Theorem~\texttt{thm:contact} and Proposition
\texttt{prop:contactcotangent}, proves the finite labelled contact
forcing and its nonzero perfect-residue value. Its complete original
LaTeX accompanies this continuation. We retain that result and compute
its receiver in the derivative of the original zero trace, with every
analytic-unit term. This is a comparison of two specified evolutions;
it does not change the definition of the arithmetic heat family.
The human heat source is Brad Rodgers and Terence Tao,
\href{https://arxiv.org/abs/1801.05914v5}{\emph{The de Bruijn--Newman
constant is non-negative}}, original equations \texttt{hoz},
\texttt{htdef}, \texttt{sas}. The original residue/trace conventions
are proved in
\href{https://github.com/KokunoYumeto/zeta-function-research-reader/blob/3c022a0adde0a6aa3d8fc8e43ef42795d88e44b0/workbenches/splitzero-tandem/continuations/20260923-supported-zero-collision-holonomy/HEAT_COLLISION_HOLONOMY.tex}{HH23--28}.

\paragraph{BT1. The specified trajectories retain the environment.}
Let $I$ be a real time interval, $x_1,\ldots,x_n:I\to\mathbb C$
locally absolutely continuous, and $U(t,z)$ a nonvanishing function
on an open spatial neighbourhood of all trajectories. Assume $U$ is
$C^1$ in real time and holomorphic in $z$, with the displayed spatial
derivatives continuous locally. Put
\[
 P=\prod_{i=1}^n(z-x_i),\quad F=UP,\quad
 g=\frac{U_z}{U},\qquad
 \iota(a)=\begin{cases}a^{-1},&a\ne0,\\0,&a=0.\end{cases}
 \tag{BT1}
\]
We compute on the exact extension with all labels retained:
\[
 \dot x_i=2\sum_{j\ne i}\iota(x_i-x_j)+2g(t,x_i)
 \quad\hbox{almost everywhere.}
 \tag{BT2}
\]
This includes the original unit drift. It is the ordinary simple-root
equation of $F_t=-F_{zz}$ when the roots are distinct and that PDE
holds, by differentiation at each root. It specifies the additional
rule at coincidence and therefore is not an assertion that the
original equation has prescribed such a rule.

Define $P_i=\prod_{j\ne i}(z-x_j)$ and
$P_{ij}=\prod_{\ell\ne i,j}(z-x_\ell)$. The complete defect is
\[
\begin{split}
 D(t,z)&=2\sum_{i<j}{\bf1}_{x_i=x_j}P_{ij},\\
 B(t,z)&=\frac{U_t+U_{zz}}{U}
       +2\sum_i\frac{g(t,z)-g(t,x_i)}{z-x_i},\\
 \boxed{F_t+F_{zz}=U D+F B.}
\end{split}\tag{BT3}
\]
Each divided difference in $B$ is holomorphic at $z=x_i$, with
value $g_z(t,x_i)$, including coincident labels.
To prove the identity, pair the reciprocal terms in
$P_t=-\sum_i\dot x_iP_i$. The polynomial identity
$P_i-P_j=(x_i-x_j)P_{ij}$ gives
\[
 P_t+P_{zz}=D-2\sum_i g(t,x_i)P_i.
 \tag{BT4}
\]
Now expand $(\partial_t+\partial_z^2)(UP)$ and use
$P_z=\sum_iP_i$. The remaining terms are
$(U_t+U_{zz})P+2U\sum_i(g(t,z)-g(t,x_i))P_i=FB$,
using $(z-x_i)P_i=P$. This proves BT3
coefficientwise almost everywhere and distributionally on compact
time intervals. All terms are locally integrable: the polynomial
coefficients and $U$ have the stated regularity, and $D$ is a finite
sum of bounded measurable coefficient functions.

\paragraph{BT2. No analytic unit cancels the contact class.}
At a fixed time let $c_1,\ldots,c_r$ be the distinct centres and
$m_1,\ldots,m_r$ their multiplicities. Then
\[
 D=\sum_{l:m_l\ge2}m_l(m_l-1)(z-c_l)^{m_l-2}
                  \prod_{h\ne l}(z-c_h)^{m_h},\qquad
 \frac DP=\sum_{l:m_l\ge2}\frac{m_l(m_l-1)}{(z-c_l)^2}.
 \tag{BT5}
\]
The second identity is meromorphic and exact, not an asymptotic
truncation. In the local coordinate $x=z-c_l$, put
$u(x)=U(t,c_l+x)\prod_{h\ne l}(c_l+x-c_h)^{m_h}$.
It is the full original local unit of $F=x^{m_l}u(x)$. Modulo
$x^{m_l}$ the defect in BT3 is
\[
 q_l=m_l(m_l-1)x^{m_l-2}u(x).
 \tag{BT6}
\]
All other contact summands and $FB$ vanish in that local algebra.
The leading coefficient in BT6 is nonzero for $m_l\ge2$.
Consequently $F_t+F_{zz}=0$ forces every collision set to have
measure zero; the analytic environment term cannot cancel contact.
The exact converse keeps the environment: the PDE holds if and only
if $D=0$ and $B=0$ almost everywhere. Indeed if the PDE holds, BT6
forces $D=0$ at almost every time; BT3 then gives $FB=0$ and hence
$B=0$ by holomorphy. Conversely both zero terms imply the PDE.
When $U=1$, $B=0$ identically and this recovers precisely the
source's residence criterion. No converse that discards $U$ is used.

\paragraph{BT3. The residue detects a derivative that the value trace can kill.}
At one centre of multiplicity $m$ retain
\[
 A=\mathbb C[x]/(x^m),\quad F_0=x^mu(x),\quad
 \lambda_u(a)=\operatorname{Res}_{x=0}\frac{a(x)}{x^mu(x)}\,dx.
 \tag{BT7}
\]
For a spatially holomorphic test $\phi$, consider the infinitesimal
coefficient perturbation $F_0+\varepsilon Q$, with
$\varepsilon^2=0$ and $Q$ holomorphic. Differentiating the fixed
root-isolating contour gives the exact zero-trace variation
\[
 \mathcal T_\phi([Q])
  =\operatorname{Res}_{0}\phi\,\partial_x(Q/F_0)\,dx
  =-\lambda_u(\phi'Q).
 \tag{BT8}
\]
It depends only on $[Q]\in A$: adding $F_0V$ adds the derivative
of a holomorphic function and contributes no residue. The contour
formula follows by differentiating $F_x/F$ on the contour, where
$F_0$ is nonzero, and integration by parts of a meromorphic
differential proves the second equality. It does not require
choosing derivatives of colliding roots.

For the contact class $q=m(m-1)x^{m-2}u$ in BT6,
\[
 \boxed{\lambda_u(xq)=m(m-1),\qquad
        \mathcal T_\phi(q)=-m(m-1)\phi''(0).}
 \tag{BT9}
\]
The full unit cancels in these particular quotients; it remains
in $q$, $F_0$, $\lambda_u$ and BT8. Since $q/F_0=m(m-1)/x^2$,
the first formula is a simple residue and the second follows by
extracting the linear coefficient of $\phi'$.

For $m\ge3$, the original cotangent complex and its duality are
\[
 [A\xrightarrow{F_0'}A\,dx],\quad
 H^{-1}=(x),\quad H^0=A/(F_0')\,dx,
 \quad\langle a,[b\,dx]\rangle=\lambda_u(ab).
 \tag{BT10}
\]
Here $(F_0')=(x^{m-1})$ since $u(0)\ne0$. The residue pairing
on $A$ is perfect: its matrix in the full power basis is
anti-triangular with nonzero anti-diagonal $1/u(0)$. The
orthogonal complement of $(F_0')$ is its annihilator, proving
the perfect induced pairing in BT10. Thus
\[
 q\in H^{-1},\qquad
 \mathcal T_\phi(q)=-\langle q,d\phi\rangle,
 \quad\operatorname{Tr}_A(M_{qa})=0\quad(a\in A).
 \tag{BT11}
\]
This is the exact connecting map from the trace-radical contact
class to a measurable trace derivative. At $m=2$ the class is
$q=2u$, its value trace is $4u(0)$, and
$F_0'q=4u(0)^2x\ne0$ in $A$. Formula BT9 still holds, but
the $H^{-1}$ assertion is not made in that case.

\paragraph{BT4. Coordinate changes and the exact test restriction.}
The variation BT8 is defined on the entire marked coefficient
space $A$. A coordinate-change direction $Q=F_0'v$ has
\[
 \mathcal T_\phi(F_0'v)=-m\phi'(0)v(0).
 \tag{BT12}
\]
This follows from $F_0'/F_0=m/x+u'/u$ by taking its residue.
Therefore BT8 descends to the quotient $A/(F_0')$ exactly when
$\phi'(0)=0$ (for $m\ge2$). When $\phi'(0)\ne0$, taking
$v=1$ exhibits the failure and computes its full correction.
The quotient has thus not been substituted for the original marked
deformation. For reflected products of tests that both vanish at
the collision, the derivative condition holds automatically; this
is the test class in which the contact pairing descends.

\paragraph{BT5. The exact arithmetic coefficient and its sign.}
Use the original arithmetic variable
\[
 s=\rho+ix/2,\quad \rho=1/2+ic_l/2,\qquad
 G(t,s)=16F(t,-2i(s-1/2)),\quad
 \mathcal E=\partial_tG-\tfrac14G_{ss}.
 \tag{BT13}
\]
The constants follow from $(-2i)^2=-4$. Consequently
$\mathcal E=16(F_t+F_{zz})$ and
\[
 \partial_t(G_s/G)
 =\tfrac14\partial_s(G_{ss}/G)+\partial_s(\mathcal E/G),
 \qquad
 \frac{\mathcal E}{G}=-\sum_l\frac{m_l(m_l-1)}{4(s-\rho_l)^2}
                  +B(t,-2i(s-1/2)).
 \tag{BT14}
\]
All terms involving the environment in the second summand are
holomorphic on the retained cluster neighbourhoods. Thus for any
fixed holomorphic arithmetic test $A(s)$, the contact contribution
to the actual derivative of the finite zero trace is
\[
 \boxed{\delta_{\rm contact}\dot Z(A)
             =\sum_l\frac{m_l(m_l-1)}4 A''(\rho_l).}
 \tag{BT15}
\]
This is BT9 with $\phi(x)=A(\rho+ix/2)$, or directly the residue
of $A\partial_s(\mathcal E/G)$ in BT14. The factor $1/4$ and
its sign are forced by the original arithmetic coordinate.

If the centre is real in the spatial coordinate, $\rho=1-\bar\rho$.
For the original involution $F^\#(s)=\overline{F(1-\bar s)}$,
take two arithmetic tests $F,G$ vanishing at $\rho$ and put
$A=F^\#G$. Direct differentiation gives
$A''(\rho)=-2\overline{F'(\rho)}G'(\rho)$. Hence
\[
 \delta_{\rm contact}\dot Z(F^\#G)
       =-\frac{m(m-1)}2\overline{F'(\rho)}G'(\rho).
 \tag{BT16}
\]
This is an evaluated negative first-jet response of the added
contact term. The genuine coefficient heat flow has the opposite
positive first-jet response. To prove this last comparison with
the complete unit, write $G_0=(s-\rho)^m v(s)$; the principal
part of $G_0''/(4G_0)$ is
\[
 \frac{m(m-1)}{4(s-\rho)^2}
      +\frac m2\frac{v'(\rho)}{v(\rho)}\frac1{s-\rho}.
 \tag{BT17}
\]
Differentiate in $s$ and take the residue against $A$. The true
heat variation is
\[
 \dot Z_{\rm heat,\rho}(A)
 =-\frac{m(m-1)}4 A''(\rho)
   -\frac m2\frac{v'(\rho)}{v(\rho)}A'(\rho).
 \tag{BT18}
\]
Adding BT15 cancels its first term exactly. The retained unit
drift remains. For the vanishing test pair, $A'(\rho)=0$ also,
so the totalized persistent-contact flow has zero derivative on
that first-jet test, whereas the genuine heat family has the
positive value. This explicitly measures the changed dynamics.

\paragraph{BT6. Fixed original Cauchy tests observe the full contact matrix.}
Let $z_a,z_b$ be distinct or equal poles with real parts greater
than one, $F_z(s)=1/(s-z)$, and
$A_{ab}=F_{z_a}^\#F_{z_b}$. Put
$\alpha=1-\bar z_a$, $\beta=z_b$, and for a centre $\rho_l$
put $d_l=\rho_l-\alpha$, $e_l=\rho_l-\beta$.
As these poles lie outside the original critical strip, they do
not meet an original arithmetic zero. In any retained local
continuation they are kept away from the cluster. Then
\[
 A_{ab}''(\rho_l)=-2\left(
 \frac1{d_l^3e_l}+\frac1{d_l^2e_l^2}+\frac1{d_le_l^3}\right),
 \tag{BT19}
\]
and the exact contact matrix is
\[
 (\delta_{\rm contact}\dot C)_{ab}
 =-\sum_l\frac{m_l(m_l-1)}2
 \left(\frac1{d_l^3e_l}+\frac1{d_l^2e_l^2}+\frac1{d_le_l^3}\right).
 \tag{BT20}
\]
Twice differentiating $A_{ab}(s)=-1/((s-\alpha)(s-\beta))$
proves BT19 and BT15 gives BT20. The original fixed-test frame
HR7 maps this matrix into the full collision algebra; no test
coefficient is discarded. On its subspace of tests vanishing at
a critical collision, the single-centre matrix reduces to the
rank-one negative form BT16. This remains a contribution to
the changed-flow derivative, not a negative value of the complete
original arithmetic form at zero.

\paragraph{BT7. Exact return to the genuine equation and complete support.}
For $U=1$ the source's finite return map is
\[
 P(t)=e^{-(t-t_0)\partial_z^2}P(t_0)
       +\int_{t_0}^t e^{-(t-r)\partial_z^2}D(r)\,dr.
 \tag{BT21}
\]
The exponential terminates on degree-$n$ polynomials, and
differentiating after the inverse exponential proves BT21
coefficientwise for absolutely continuous paths. For the actual
unit the complete returned equation is BT3, including $B$;
there is no omitted external-root contribution.

Finally the input inverse is exactly the generalized inverse in
the original $G_L(\mathbb C)$. In its top supported ring fibre,
\[
 d(a)=1-aa^\dagger,
 \qquad p(d(a))={\bf1}_{p(a)=0},\qquad
 p(D)=2\sum_{i<j}p(d(x_i-x_j))P_{ij}.
 \tag{BT22}
\]
Here differences are taken in the top supported fibre, whose
additive identity is $e$. Thus the contact indicator detects
supported equality and has never identified it with absence
$\tau$. The linear map $Q\mapsto\mathcal T_\phi(Q)$ lifts as a
$G_L(\mathbb C)$-semimodule map on the full supported carriers:
$(Q,1_L)\mapsto(\mathcal T_\phi(Q),1_L)$ and
$(0,\lambda)\mapsto(0,\lambda)$. Linearity of BT8 preserves
amplitude addition and scalar action; unchanged labels preserve
support joins and scalar meets. The residue pairings lift as
bilinear maps, with output support the meet of the input supports.
Distributivity of $L$ verifies both additive laws, and residue
bilinearity verifies the amplitudes. Zero amplitudes retain their
declared support labels. In particular the
zero value trace of BT11 remains a supported zero, while BT9
gives a nonzero tangent observation on the same nonzero class.

\begin{figure}[p]
\centering
\begin{tikzpicture}[>=Stealth,box/.style={draw,rounded corners,
align=center,text width=11cm,inner sep=9pt}]
\node[box] (q) at (0,0) {Full contact class $q=m(m-1)x^{m-2}u$\\
$F_0=x^mu$; all original unit coefficients and support labels retained};
\node[box] (r) at (0,-2.2) {Perfect residue: $\lambda_u(xq)=m(m-1)$\\
For $m\ge3$: $q\in H^{-1}$ while $\operatorname{Tr}M_{qa}=0$};
\node[box] (t) at (0,-4.4) {Derivative receiver $\mathcal T_\phi(q)=-\lambda_u(\phi'q)$\\
Original coordinate $s=\rho+ix/2$: $\delta\dot Z(A)=m(m-1)A''(\rho)/4$};
\node[box] (v) at (0,-6.8) {Tests vanishing at a critical collision\\
contact: $-m(m-1)\overline{F'(\rho)}G'(\rho)/2$\\
genuine heat: $+m(m-1)\overline{F'(\rho)}G'(\rho)/2$};
\draw[->] (q)--node[right]{BT7--11}(r);
\draw[->] (r)--node[right]{cotangent pairing with $d\phi$}(t);
\draw[->] (t)--node[right]{BT16--18: exact cancellation}(v);
\end{tikzpicture}
\caption{The contact class has a nonzero trace-derivative receiver even
when its ordinary trace pairing vanishes. Its added derivative cancels
the genuine heat splitting term. These are two computed contributions
to the specified equations; their difference retains the original clock,
analytic unit and arithmetic tests. Complete proofs BT1--22; the
supplied contact derivation and human heat source are cited above.}
\end{figure}
